\documentclass[12pt]{jlreq}
\usepackage{amsmath,amssymb,amsthm}
\theoremstyle{definition}
\newtheorem*{solution}{解答}

\begin{document}

\begin{solution}
これは $\ker \rho_1 \cap \ker \rho_2$ に含まれる。

H24 B4
(1) $e_2 = efe$ とし、$f = ef + fe$ に $xe$ をかけると、
$fe = efe + fe^2$
$efe = 0$
$a_2(1-e) = (1-e)f(1-e)$ とし、$f = (1-e)f + f(1-e)$。
これに $x(1-e)$ をかけると、
$f(1-e) = (1-e)f(1-e) + f(1-e)^2$
$(1-e)f(1-e) = 0$

(2) $e^2 = e$ in $B$ とし、$\varphi: B \to \mathbb{C}$ 1次元 rep とすると、$\varphi(e) = 0$ or $1$。
$efe = 0, (1-e)f(1-e) = 0$ なので、どちらに進む $\varphi(f) = 0$。
$A \to \mathbb{C}$
$e \mapsto 0$ の $\ker \supset I$
$f \mapsto 0$
$A \to \mathbb{C}$
$e \mapsto 1$ の $\ker \supset I$
$f \mapsto 0$
より、$\varphi_0, \varphi_1: B \to \mathbb{C}$
$e \mapsto 0$ or $1$
$f \mapsto 0$
が $B$ の次元 rep。

(3) $\varphi: B \to GL(V)$ rep とする ($\dim V \ge 2$)。
$V_1 = \ker \varphi(1-e)$
$V_2 = \ker \varphi(e)$
とすると $V = V_1 \oplus V_2$。
$v \in V, v = \underbrace{ev}_{V_1} + \underbrace{(1-e)v}_{V_2}$
$v \in V_1 \cap V_2$ とすると $v = 1 \cdot v = e \cdot v + (1-e) \cdot v = 0$。
$\varphi(e)|_{V_1} = \text{id}_{V_1} \to \varphi(e) = \text{id}_{V_1} \oplus 0_{V_2}$
$v \in V_1$ とすると、$(1-e)v = 0$ より $v = ev$。
同様に $\varphi(1-e) = 0_{V_1} \oplus \text{id}_{V_2}$
$\varphi(f)(V_1) \subset V_2, \varphi(f)(V_2) \subset V_1$
$v \in V_1, e(f \cdot v) = ef \cdot v = ef \cdot (ev) = efe \cdot v = 0$。
$\ker (\varphi f)|_{V_1} \neq 0$ なら、$V$ の sub rep として1次元表現がある。
よって $\varphi(f)|_{V_1}: V_1 \to V_2$ は単射である。
同様に $\varphi(f)|_{V_2}: V_2 \to V_1$ は単射。
$\varphi(f^2)|_{V_1}: V_1 \to V_1$ の固有ベクトル $v$ をとる。
$\mathbb{C}v \oplus \mathbb{C}fv$ は $e$ と $f$ の作用で不変（ただし $f \cdot v \neq 0$）。
よって $V: \text{既約} \implies \dim V = 2$。
$v \in V_1, v \neq 0$ をとり（これに固有な $f^2$ の固有値になるので）、$f^2$ の固有値を $a \in \mathbb{C}$ とする。
このとき $v, f \cdot v$ に関する $e, f$ の表現行列は
$e = \begin{pmatrix} 1 & 0 \\ 0 & 0 \end{pmatrix}, f = \begin{pmatrix} 0 & a \\ 1 & 0 \end{pmatrix}$
$A \to GL(2, \mathbb{C})$ に
$e \mapsto \begin{pmatrix} 1 & 0 \\ 0 & 0 \end{pmatrix}$
$f \mapsto \begin{pmatrix} 0 & a \\ 1 & 0 \end{pmatrix}$
とすると、$GL(2, \mathbb{C})$ で $e^2 = e$。
$ef + fe = \begin{pmatrix} 0 & a \\ 0 & 0 \end{pmatrix} + \begin{pmatrix} 0 & 0 \\ 1 & 0 \end{pmatrix} = f$
となり、$B \to GL(2, \mathbb{C})$ が rep となる。
（ちなみに $f$ の特性多項式は $x^2 - a$ で $a$ が変えようのでこれらは全て異なる）
キャラクタのことは $f: V_1 \to V_2, V_2 \to V_1$ がひっくり返って折り返る。

Def. 体拡大 $L/K$ が分離であるとは、
$\forall \alpha \in L, \alpha$ の $K$ での最小多項式が分離であること。（$f'$ の因子を持たない）
$L/K$ が純非分離拡大であるとは、
$\forall \alpha \in L, \alpha$ の $K$ 上の最小多項式が純非分離。

Prop.
$L/K$: 先非分離 $\iff \forall \alpha \in L, \exists n \ge 0, \alpha^{p^n} \in K$
\end{solution}

\end{document}